\chapter{Introduction}



Particle physics is dedicated to unraveling the most fundamental constituents of the universe and the forces that govern their interactions. It explores the known elementary particles, such as quarks, leptons, and bosons, which constitute matter and mediate the fundamental forces --- gravity, electromagnetism, the weak nuclear force, and the strong nuclear force. The Standard Model is the central theoretical framework in this field, describing these particles and forces. A theory that was highly successful in explaining the behavior of most known particles and forces but which initially assumed that neutrinos were massless faced challenges when experimental evidence showed that neutrinos do have mass and undergo flavor-changing oscillation. This discovery revealed the need for a modification in this model to give a complete description of the universe’s fundamental particles and forces. As a result, pursuing a more comprehensive theory capable of explaining neutrino mass and other phenomena beyond the Standard Model became a key focus in the field. Therefore, studying neutrinos and their interactions is crucial because it opens doors to exploring new theoretical models and new physics.


Neutrinos were theoretically proposed by Pauli in 1930 to preserve angular momentum and energy conservation in beta decay. They were experimentally discovered in 1956 by Cowan and Reines after many years of not having the necessary technology for their detection due to their limited interaction with matter. Since then, neutrinos have been extensively studied to gain a deeper understanding of them, their properties, the way they interact, whether they are Dirac or Majorana neutrinos, their masses, etc. Over the years, these studies have led to advancements in both theoretical and experimental domains, as the limited interaction of these particles has demanded new technologies and improvements in aspects such as detector sizes and time scales.

This work focuses on Coherent Elastic Neutrino-Nucleus Scattering
(CE$\nu$NS), a process whose cross-section is accurately defined by
the Standard Model, so any deviation from its expected outcome could
indicate the presence of new physics phenomena, making it a helpful
process for exploring new physics. However, there are several
challenges that this type of measurement faces. First, the observation
itself is a great challenge due to the low momentum transfer, which
implies low-energy nuclear recoils. Besides, in some cases, there are
nuclear effects that need to be taken into account. In this thesis, we
study the expectations for a combined measurement of CE$\nu$NS using
two different detectors made from isotopes of the same element. 
We explore the advantage of using isotopes
with different numbers of neutrons in this simultaneos detection.
A
similar idea was put forward recently~\cite{Galindo} using germanium
detectors as an example.
We
will show how these experimental setups may help to improve the
sensitivity to new physics, using as an example the formalism of
non-standard interactions (NSI).


The thesis introduces, in Chapter 2, a brief introduction to the
Standard Model and the phenomenology of neutrinos, including its
interactions and the vacuum oscillation phenomenon. In Chapter 3, we
focus on the CE$\nu$NS interaction, introducing its cross-section and
the corresponding nuclear form factors that can be used. We will also
describe the possible neutrino sources and present the computations of
the expected number of events for the isotopes that we are interested
in. Chapter 4 describes our proposed isotopes for being used as
a CE$\nu$NS detector and shows the expected sensitivity to standard physics and  NSI. Finally, we give our conclusions in Chapter 5. 





